\chapter{Primeros Pasos con Python desde Lisp}
\label{chap:python-firststeps}
{\huge$\bigstar$}{\huge$\bigstar$} Intermedio

\section{Verificación}
\begin{lispcode}
(find-package :py)
\end{lispcode}

\section{Importar Módulos}
\begin{lispcode}
(let ((math (py:import "math"))) (py:call (py:attr math "sqrt") 16))
\end{lispcode}

\section{Builtins Directos}
\begin{lispcode}
(py:print "Hola")
(py:abs -42)
(py:sum '(1 2 3 4 5))
(py:sorted '(3 1 4 1 5))
(py:len "hello")
\end{lispcode}

\section{Caché de Módulos}
\begin{lispcode}
(time (py:import "math"))  ;; segunda vez: instantáneo
\end{lispcode}

\section{Llamada Directa}
\begin{lispcode}
(let ((sqrt (py:attr (py:import "math") "sqrt"))) (sqrt 144))
\end{lispcode}

\begin{examplebox}[title=Usando math desde Lisp]
\begin{lispcode}
(let ((math (py:import "math")))
  (format t "pi = ~F~%" (py:attr math "pi"))
  (format t "sqrt(16) = ~F~%" (py:call (py:attr math "sqrt") 16)))
\end{lispcode}
\end{examplebox}

